\section{Conclusion}
\label{sec:conclusion}
This paper presents a reproducible diagnostic protocol for sparse-concept and calibration evaluation in Knowledge Tracing. The protocol combines train-only KC-frequency stratification, learner-based and temporal splits, limited cold-start concept analysis, ECE and Brier decomposition, reliability diagrams, and an eight-channel leakage and predictive sanity audit. The resulting artifact is intended to support fairer and more informative evaluation of KT models, especially future structure-aware and self-supervised approaches targeting sparse KCs. For instance, our experimental application of the protocol reveals that while aggregate population metrics may remain competitive, model calibration can degrade significantly on sparse concepts---such as SimpleKT's Expected Calibration Error (ECE) doubling from $0.11$ on dense concepts to $0.22$ on sparse concepts in ASSISTments 2012. More broadly, our temporal-split audit demonstrates that the protocol's built-in leakage and sanity checks successfully catch subtle evaluation pipeline bugs, ensuring that researchers can distinguish genuine concept-level cold-start failure from alignment artifacts.
